\documentclass[UTF8,a4paper,11pt]{ctexart}
\usepackage[margin=24mm,headheight=15pt]{geometry}
\usepackage{booktabs,longtable,array,amsmath,hyperref,fancyhdr}
\hypersetup{hidelinks}
\setlength{\parindent}{2em}
\setlength{\parskip}{0.25em}
\renewcommand{\baselinestretch}{1.23}
\pagestyle{fancy}
\fancyhf{}
\fancyhead[L]{“韬”智能杯 EDA 创新大赛 · 赛题二}
\fancyhead[R]{AIC-2026-18059446}
\fancyfoot[C]{\thepage}
\renewcommand{\headrulewidth}{0.3pt}
\ctexset{section={format=\large\bfseries,beforeskip=1.1em,afterskip=0.4em}}
\begin{document}
\begin{center}
{\LARGE\bfseries 面向 3D 芯片网表切分的优化目标演化}\par
\vspace{0.35em}
{\Large 技术方案}\par
\vspace{0.6em}
参赛编号：AIC-2026-18059446\quad 所选基线：概率预测 GNN（\texttt{binary\_loss}）\par
2026 年 11 月 20 日
\end{center}

\section{算法概述}
赛题将网表实例分配到双层芯片的上、下层，切分结果通过固定的布局、供电网络、层间端子与时序优化流程评估。本方案选择概率预测 GNN 基线，只修改训练损失函数。GNN 结构、训练设置、切分阈值、约束修复及下游物理流程均沿用官方任务包。

我们实现的 Agent 接受官方初始程序、调用统一 DeepSeek-V4.1-Flash 模型生成候选、运行官方静态检查、提交完整训练矩阵、读取时序反馈，并保留已完成评估中得分最高的候选。五次本地候选提交中有两次完成完整矩阵；选定候选的三个训练算例、每例三个种子全部通过核验，训练集得分为 28.972451。该值只对应本地训练集，不是七算例正式成绩。


传统网表切分常以割网数量和层间面积均衡为近端目标，后续布局和时序分析才给出最终物理效果。赛题的固定基线以 GNN 预测每个实例的上层概率，损失中同时含期望割网、软面积偏差和二值化正则，再由固定阈值及修复程序形成合法切分。近端损失与最终 WNS、TNS 并不完全一致，因此需要把完整物理流程的反馈返回给 Agent，而不能只根据训练损失选择候选。

本方案的改动幅度保持较小：最终选定的损失函数只改变软面积项的形状。这样可以把观察到的时序变化与一项明确的代码差异对应起来，也便于检查是否仍满足官方函数接口。对于最终成绩，我们只引用官方评测器对候选程序生成的原生物理结果；训练结果不外推到未运行的算例。

赛题的困难在于，切分器训练时看到的是连续概率和可微代理损失；物理设计流程使用离散的上下层归属，并进一步进行面积修复、布局和时序优化。软面积比接近目标只表示训练目标中的一项较小，不能保证修复次数少，更不能保证关键路径改善。本方案因此把所有候选放进同一原生流程，以时序反馈决定是否保留，而不是用损失下降或单个切分统计量代替比赛指标。
\section{实现方案}
\textbf{Agent 工作流程。} 提交目录的 \texttt{agent.py} 实现 \texttt{Agent.run(context)}。每轮将初始程序或当前已核验的最佳程序、简短的运行记忆和最新反馈送入官方模型代理；从回复中提取完整函数；先调用任务包静态校验器，失败时允许一次修复；检查是否与父代相同，再提交训练矩阵并等待原生物理评估。每次评估结束后记录候选、反馈和失败原因，只把完整评分结果用于最优候选选择。正式提交的版本可在官方预算内自动继续迭代，返回已完成训练评估的候选 ID。模型调用、评估次数及物理流程由官方外层框架控制。


\textbf{提示词与候选控制。} 初始提示词先给出允许修改的函数范围、可用数值接口及禁止的操作，再给出不超过官方篇幅限制的通用离线知识。每轮用户侧提示包含当前父代程序、已完成评估的简要历史、上一轮反馈和剩余评估额度。Agent 要求模型一次只改变一个机制，并返回完整函数文本。生成代码的函数名、参数顺序、语法和运算范围由官方静态校验器检查；静态检查不消耗候选评估次数。修复仍失败或生成结果与父代等价时，Agent 放弃该草案，不将其当作已评分候选。

\textbf{评估与恢复。} 正式候选的一个评估提交覆盖三个训练算例、每例三个种子。Agent 读取九组的分数、WNS、TNS、切分面积和失败状态，记录每次提交的代码与父代选择。评估中断后，程序依据外层框架保存的请求回执恢复状态；只有取得完整训练矩阵和候选源码的成功记录才可以更新最优候选。已经受理而状态不明的请求先对账，避免重复提交。最终版本把每次会话的提交上限设为官方预算允许的 10 次；费用和次数仍由组委会外层框架约束。此前本地逐次运行时，每次会话只提交一个新候选，便于检查结果；这与最终自动连续运行的入口有所区别。
\textbf{损失函数。} 设 $p_i$ 为实例 $i$ 归属上层的概率，$a_i$ 为其面积，$r=\sum_i a_i p_i/\sum_i a_i$ 为软上层面积比例，$r_0$ 为官方目标。初始损失由期望割网、$(r-r_0)^2$ 面积项和概率二值化项组成。本地选定程序保留割网项与二值化项，将面积项改为
\[
L_{\mathrm{area}}=\bigl[\max(|r-r_0|-0.03,\,0)\bigr]^2 .
\]
该 0.03 死区允许软面积比例在目标附近小幅浮动，超出后继续施加二次惩罚；切分后的面积与宏块合法性仍由官方修复与检查决定。函数未改变调用签名、返回值类型、GNN、优化器或物理评估逻辑。其作用是减少过强软面积约束对割网与二值化目标的干扰；实际收益仍须按原生时序结果判断。


二值化项保持初始程序的形式。记 $m_i=[\max(0.9-|2p_i-1|,0)]^2/0.81$，则
\[
L_{\mathrm{bin}}=0.25\,\mathrm{mean}_i(m_i)
+0.75\frac{\sum_i a_i m_i}{\sum_i a_i}.
\]
完整目标仍为官方传入权重乘以期望割网、面积项和二值化项之和。面积死区只作用于训练时的软概率，不改变面积平衡的最终合法范围，也不覆写官方的修复过程。候选程序不读取算例名、种子或隐藏数据；所有算例共用同一个函数。

\textbf{接口边界。} 提交的算法文件只实现任务包约定的损失函数，输入包括实例概率、面积、超边信息、目标面积比及三个已有权重，输出为可反向传播的标量。面积比计算使用所有实例面积之和作为分母；二值化项同时采用节点均值和按面积加权的均值，使小实例和大实例都受到约束。代码中没有读取文件、环境变量或评测目录的分支，也没有按算例名称设置系数。Agent 的生成、检查和选择过程写在独立入口中；获选的程序快照单独放在同目录，供初筛复跑。
\section{实验与创新点}
\textbf{设置与得分。} 本地使用官方 \texttt{contest-source.tar.gz}、\texttt{contest-data.tar.gz} 与固定运行镜像，在 \texttt{bp\_fe-ng45}、\texttt{bp-ng45}、\texttt{ariane133-ng45} 三个训练算例上各运行种子 1001、1002、1003。时序边界为 OPT 后、CTS 前。表 1 将任务包提供的初始程序参考结果与选定候选的原生 WNS/TNS 逐组列出；最后一列是评估器按共同参考值计算的得分，不能直接由表中两组数值相减得到。

\begin{center}\scriptsize
\begin{tabular}{llrrrrr}
\toprule
算例 & 种子 & 初始 WNS & 候选 WNS & 初始 TNS & 候选 TNS & 候选得分\\
 & & \multicolumn{2}{c}{ns} & \multicolumn{2}{c}{ns} & \\
\midrule
bp\_fe-ng45 & 1001 & -0.6980 & -0.7861 & -506.96 & -737.38 & 3.359\\
 & 1002 & -0.8543 & -0.8487 & -313.49 & -404.00 & 28.481\\
 & 1003 & -0.8388 & -0.7020 & -604.36 & -138.18 & 57.242\\
bp-ng45 & 1001 & -8.2449 & -9.7678 & -5654.09 & -7763.76 & -15.436\\
 & 1002 & -8.3899 & -8.0861 & -8888.10 & -7856.77 & -8.514\\
 & 1003 & -8.4867 & -8.2367 & -9086.50 & -6300.17 & 4.893\\
ariane133-ng45 & 1001 & -0.6537 & -0.2794 & -15.53 & -1.94 & 77.673\\
 & 1002 & -0.7316 & -0.1658 & -210.98 & -0.56 & 87.117\\
 & 1003 & -0.1847 & -0.4944 & -0.20 & -58.95 & 25.938\\
\bottomrule
\end{tabular}\par
\smallskip
表 1：初始程序与选定候选的训练集原生时序指标（负值，越接近 0 越好）。
\end{center}


\textbf{演化记录。} 表 2 将五次正式候选提交按时间列出。括号中的短编号对应 \texttt{logs/} 中的完整候选 ID 和代码快照。父代列表示该次生成前 Agent 所选的当前最优程序；失败或较低分不会替代已核验的最佳父代。
\begin{center}\small
\begin{tabular}{cllrl}
\toprule
轮次 & 候选短编号 & 父代 & 训练得分 & 结果\\
\midrule
1 & 64a1b258 & 官方初始程序 & 28.972 & 完整 9/9，选定\\
2 & 91b87b3b & 第 1 轮 & -- & 运行中断，矩阵不完整\\
3 & f3dac565 & 第 1 轮 & 7.250 & 完整 9/9，未选定\\
4 & e495a6c7 & 第 1 轮 & -- & 运行中断，矩阵不完整\\
5 & 1ee0c61e & 第 1 轮 & -- & 部分组合内存不足\\
\bottomrule
\end{tabular}\par
\smallskip 表 2：候选关系与本地训练评估状态。
\end{center}

\textbf{切分诊断。} 为了检查面积死区是否真正改变后续切分，表 3 汇总初始程序与选定候选的三个种子平均值。原始上层面积比例来自修复之前的概率切分；“修复搬移面积”表示官方平衡修复所搬移的面积占比。两者是切分诊断，不是最终时序得分。
\begin{center}\small
\begin{tabular}{lrrrr}
\toprule
算例 & 初始原始上层比 & 候选原始上层比 & 初始修复搬移比 & 候选修复搬移比\\
\midrule
bp\_fe-ng45 & 0.513 & 0.499 & 0.130 & 0.066\\
bp-ng45 & 0.391 & 0.464 & 0.413 & 0.413\\
ariane133-ng45 & 0.518 & 0.534 & 0.220 & 0.355\\
\bottomrule
\end{tabular}\par
\smallskip 表 3：按算例对三个种子取均值的切分诊断。
\end{center}
为区分“面积更平衡”与“连接更容易布局”，还汇总了切网比例和时序加权割网代价。表 4 同样先对每个算例的三个种子求均值；数值来自评估产物的切分诊断字段，不是评分公式中的独立加分项。切网比例越低通常意味着层间连接较少，但最终时序仍取决于网络位置、布线以及关键路径。
\begin{center}\small
\begin{tabular}{lrrrr}
\toprule
算例 & 初始切网比例 & 候选切网比例 & 初始时序加权割网 & 候选时序加权割网\\
\midrule
bp\_fe-ng45 & 0.10987 & 0.11738 & 5419.17 & 5146.42\\
bp-ng45 & 0.10202 & 0.10265 & 16844.83 & 16970.91\\
ariane133-ng45 & 0.08596 & 0.08079 & 22869.05 & 20933.02\\
\bottomrule
\end{tabular}\par
\smallskip 表 4：三个种子的平均连接代价诊断；两个指标的定义不同。
\end{center}
在 \texttt{bp\_fe-ng45} 中，候选的普通切网比例略升，而时序加权割网下降；在 \texttt{ariane133-ng45} 中两项均下降；\texttt{bp-ng45} 中两项均略升。这与各案例的时序收益方向大致一致，但不能单凭相关性推断具体关键路径如何变化，也不能替代 OPT 后的时序报告。训练目标中的期望割网是在连续概率上计算的，表 4 则是最终切分的离散统计，两者含义也不相同。
三个算例的候选平均得分依次为 29.694、-6.353、63.576，等权平均为 28.972。\texttt{bp-ng45} 未改善，说明单一面积死区机制并非对所有电路有效。训练运行有明显种子波动，尤其 \texttt{ariane133-ng45} 的种子 1003。五次提交共消耗 5/10 次评估额度和约 0.481 元模型费用；另一个完整候选得分为 7.250，未超过选定版本。其余三次矩阵未完整完成，未用于候选选择。本地未运行公开验证，故此处不填写公开验证数值；正式成绩以组委会统一评测为准。

\textbf{结果判读。} 对于 \texttt{bp\_fe-ng45}，候选的原始面积比例更接近 0.5，修复搬移面积均值约减半；其中种子 1003 的时序改善最大，而种子 1001 的 WNS、TNS 都比初始参考差。\texttt{bp-ng45} 的原始面积比例从约 0.391 提高到 0.464，但修复搬移比例基本不变，三个种子中两个得分为负。\texttt{ariane133-ng45} 的三个种子均得到正分，不过修复搬移面积增加，说明面积修复量与时序结果也不是简单的单调关系。这些观察只描述同一固定流程下的结果，不足以断言死区改善了所有电路的布线或时序。

表 1 的初始程序列来自任务包提供的参考记录，候选列来自所选请求的评估结果；每一行均保留相应的原生日志、切分结果和指标文件。候选得分使用评估器的共同评分参考，按种子汇总后再对算例等权汇总。训练集只有三个算例、九个种子；官方正式评估涉及七个算例，因此 28.972 只能作为本地选择依据。公开验证没有完成，本报告不以估算值填补该部分。

\section{问题与思考}
本地结果表明：训练集整体得分为正，但 \texttt{bp-ng45} 平均退化，跨算例收益不均衡。种子间差异也较大，不能把最高单组得分视为稳定收益。死区可能允许概率在目标面积附近发生更大变化；这种自由度既可能减少无用的面积惩罚，也可能改变离散切分、面积修复和最终路径。固定一个超参数并不能保证每种拓扑都受益。

\textbf{失败与中断。} 三次未完成的评估中，两次与运行中断后的矩阵不完整有关，一次在部分组合上出现主机可用内存不足。它们在预算账本中依然有提交记录，但没有完整九组结果，因此没有进入最优候选比较。代码恢复逻辑先读取请求回执与本地历史，避免把正在对账的请求误当成新评估。保留这些失败记录，是为了让评委区分算法本身的物理结果与执行环境造成的缺失。

\textbf{代理目标与最终目标的差距。} 表 3 显示原始面积比例改善并不必然降低修复搬移，修复搬移减少也不必然改善每个种子的时序。期望割网只衡量训练阶段的连接代价，真实时序还受到布线位置、端子分配、优化器决策和关键路径的影响。后续可以尝试分阶段调节割网与面积权重，或把更接近最终时序的可微代理加入目标；这些方向在本次提交中没有评测证据，故未纳入最终程序。

\textbf{证据边界。} 本地采用官方默认存储模式完成最佳候选评估，保留了切分、阶段指标和时序报告；该模式会清理中间布局数据库。为补全布局文件而单独运行同一程序时，\texttt{ariane133-ng45} 的种子 1001 再次得到合法切分和物理结果，但离散切分与首次评分运行不相同，WNS/TNS 也有差异。因此补采结果与原评分回执分目录保存，不把二者拼成同一次运行，也不重新计算训练得分。公开验证与隐藏算例的泛化表现未知；最终是否通过组委会复现及取得排名，以官方复跑为准。

\section{过程进度与团队分工}
工作首先对齐官方接口、可修改范围和任务包版本，完成 Agent 的候选生成、静态检查、反馈读取与结果恢复。随后在固定运行环境中执行五次候选提交，并依据完整九组训练矩阵选择首轮候选。最后核对选定程序的哈希、原生结果、失败记录与提交目录，整理技术报告和复现入口。整个过程中只改变允许提交的损失函数和 Agent 实现，物理流程保持任务包版本。

\textbf{版本与预算。} 初始程序来自官方任务包；首个候选 64a1b258 在九组训练上取得 28.972。第三个候选 f3dac565 完成相同矩阵，取得 7.250，因而没有替换首个候选。第二、四、五个候选没有完整矩阵，保留为过程记录。提交材料将首个候选的完整源码固定为 \texttt{program.py}，并把五次提议的代码快照、请求编号和结果保存在 \texttt{logs/}。预算记录显示已使用五次本地评估；提交时的 Agent 会依据官方剩余额度继续搜索，而不是依赖本地文件中的个人密钥。

参赛团队负责方案选择、代码审查、原生运行和材料核验。DeepSeek-V4.1-Flash 通过官方代理辅助生成候选函数；外部开发工具用于编码、日志检查和报告整理。报告不把模型生成过程写成未发生的人工实验，也不把未完成的候选计为有效成绩。

\section{复现与参考资源}
复现入口、任务包校验值、运行镜像摘要及选定候选编号见同包 README。复现者先获取官方源码、数据与运行镜像，并核对版本；然后将本包 \texttt{AgenticBench/agent/binary\_loss/} 复制至任务包对应目录，安装 \texttt{requirements.txt} 中的额外依赖。先运行官方 \texttt{contest2 check} 静态检查，再以新工作目录运行 \texttt{contest2 run}，由外层框架执行候选生成、训练矩阵和七算例复评。模型密钥由评测环境提供，不在提交包内。

提交包中的训练结果目录按算例和种子保存最佳候选的原生记录；初始参考目录保存任务包提供的初始程序结果；逐种子 CSV 是两者的索引。日志目录内有候选快照、父子关系、请求反馈和预算账本。评委可先按候选 ID 与程序哈希找到所选算法，再用九组结果核对表 1，最后运行 Agent 检查其能否从初始程序自主搜索。公开验证未运行，因此材料中没有相应结果目录或成绩声明。

为避免复现时混入不同版本，提交 README 列出官方源码归档的 SHA-256、任务包记录的运行镜像摘要和本地程序哈希。物理文件以生成它们的程序、算例和种子单独标记；默认清理造成的缺失也在同一目录说明。所有实验结果都在本机按原生流程产生，最终七算例分数仍需要组委会统一运行确认。

\textbf{材料索引与核对。} 表 5 给出最直接的检查顺序。代码和物理结果分别归档，避免将生成某组结果的算法误认成另一轮候选。每个独立补采的物理结果另有清单，记录程序哈希、运行镜像、算例、种子以及输出文件的 SHA-256。补采记录只用于保存默认清理模式没有留下的布局文件；训练得分仍引用先前完整评估的回执，不从补采日志重新构造分数。
\begin{center}\small
\begin{tabular}{p{0.25\linewidth}p{0.63\linewidth}}
\toprule
材料 & 可复核内容\\
\midrule
Agent 开发目录 & 入口、提示词、静态检查、候选选择和恢复逻辑\\
最终程序 & 选定的同一损失函数及其 SHA-256\\
训练结果目录 & 九组切分、指标、时序报告及可用的物理文件\\
演化日志 & 五次候选的源码、父子关系、评估状态与失败原因\\
\bottomrule
\end{tabular}\par
\smallskip 表 5：提交包中的主要材料及其用途。
\end{center}

对于一个独立补采目录，先检查候选程序哈希是否与最终程序一致，再检查切分阶段和 OPT 阶段均显示成功；最后对照最终 WNS/TNS 与原训练评估中同一算例、同一种子的数值。上下层布局数据库和 DEF 位于该目录的 OPT 阶段，装配时序报告位于同阶段的 assembly 子目录。这些步骤能帮助评委快速定位物理结果，也明确了补采与正式评分的证据边界。

方法参考官方赛题说明和任务包，及 TP-GNN（Lu 等，DAC 2020）、Open3DBench（Shi 等，arXiv:2503.12946v2）、DREAMPlace（Lin 等，DAC 2019）与 OpenROAD 文档。OpenROAD 在官方源码包中附 BSD 3-Clause 许可证；本提交未修改或重分发其源码。物理流程及其软件版本由官方运行镜像摘要固定。模型名称、任务包来源、使用位置和版本信息均在 \texttt{README.md} 中列明。
\end{document}
